\begin{abstract}
A meeting interpreter must say what each person said, attribute it to the right person, and decide when to say it, also
when several people talk at once. In \sys{}, a speech LLM reads acoustic tokens interleaved with tokens of a frozen
streaming diarizer, writes speaker-first, time-stamped translations for every speaker of an overlapping window, and hands
them to a speech generator prompted with each speaker's voice. On AMI reference windows \sys{}
beats the strongest non-oracle diarize-then-translate cascade by \num{v7:gainoff:en2zh} (English$\rightarrow$Chinese) and
\num{v7:gainoff:zh2en} (Chinese$\rightarrow$English) speaker-attributed BLEU, and on overlapped speech it leads a
speaker-blind model relabelled by the same diarizer by \num{v8:x:ov-O2-rel-en2zh:diff} and
\num{v8:x:ov-O2-rel-zh2en:diff} BLEU. Trained on fixed per-chunk deadlines, the decoder writes much of its output only after an
utterance has ended, and a rule-based speech layer queues it. We learn both timing decisions as latent structure with
exact EM. Write alignment decides when, and in which speaker order, each translated chunk is written; one constant write
bias, set on development meetings, corrects the drift of EM towards later writes. Playback alignment decides which
rendering of each block is spoken given the playback backlog, under a posterior-regularised latency budget. Against the
same jointly trained deadline model on the two AMI test directions (Zh$\rightarrow$En / En$\rightarrow$Zh), \sys{} raises
text BLEU by \num{v10:sig:test:ami:text:m-b:diff} / \num{v10:sig:test:en:text:m-b:diff} and speech ASR-BLEU by
\num{v10:sig:test:ami:speech:m-b:diff} / \num{v10:sig:test:en:speech:m-b:diff}, and changes per-stream speech LAAL by
\num{v10:test:ami:m-b:slaal} / \num{v10:test:en:m-b:slaal}\,s. A decomposition shows the two parts are complementary:
playback alignment brings the speech gains and the whole latency reduction, write alignment the text gain.
\ifAliwin
On real Mandarin meetings (AliMeeting) it changes text BLEU by \num{v10:test:ali:m-b:all} and speech ASR-BLEU by
\num{v10:test:ali:m-b:asr}.
\else
On real Mandarin meetings (AliMeeting), held out from all training and tuning, text BLEU changes by
\num{v10:sig:test:ali:text:m-b:diff} (not significant), so we claim no text gain there, while speech ASR-BLEU (\num{v10:sig:test:ali:speech:m-b:diff})
and speech LAAL (\num{v10:test:ali:m-b:slaal}\,s) still improve.
\fi
Per utterance, the system remains closer to consecutive than to simultaneous interpretation.
\end{abstract}
